% Copyright (c) 2026 Geovana Neves. Licensed under CC BY 4.0 (https://creativecommons.org/licenses/by/4.0/). See LICENSE-AND-AUTHORSHIP.md.
%
% shared/tikz/guide-inputs-into-run.tex
% Which input file supplies which part of a run. The reader should be able to
% look at this and know which file to open.
% Shared fragment: no sectioning, no frame, no size command, no float. Wrapped by the caller.
\begin{tikzpicture}[font=\scriptsize, node distance=3mm,
  f/.style={draw=CourseBlue, thick, rounded corners, fill=CourseBlue!8,
            align=left, text width=34mm, minimum height=7mm, inner sep=3pt},
  ctr/.style={draw=CourseGold, very thick, rounded corners, fill=CourseGold!15,
              align=center, text width=26mm, minimum height=14mm, inner sep=3pt},
  ans/.style={font=\tiny\itshape, text=CourseGray, align=left, text width=32mm}]

\node[ctr] (row) {\textbf{one row of}\\ \texttt{matriz.fs}\\ \tiny cites by name};

\node[f, above left=14mm and 16mm of row] (geo) {\texttt{geometries/<mesh>/}};
\node[f, below=2mm of geo] (ref) {\texttt{references/r001.toml}};
\node[f, below=2mm of ref] (set) {\texttt{setups/s001.toml}};
\node[f, below=2mm of set] (pp)  {\texttt{pproc/p001.toml}};
\node[f, below=2mm of pp]  (exe) {\texttt{executables.toml}};

\node[ans, left=2mm of geo] {what shape};
\node[ans, left=2mm of ref] {what it is divided by};
\node[ans, left=2mm of set] {how hard to solve};
\node[ans, left=2mm of pp]  {what to export};
\node[ans, left=2mm of exe] {which build};

\foreach \n in {geo, ref, set, pp, exe}
  \draw[-{Stealth}, thick, CourseBlue] (\n.east) -- (row.west);

\node[f, right=16mm of row, text width=30mm, draw=CourseGray, fill=CourseGray!8]
  (emit) {the emitted script\\ \tiny one per point};
\draw[-{Stealth}, thick, CourseGold] (row) -- (emit);

\node[font=\tiny\itshape, text=CourseGray, below=6mm of row, text width=80mm, align=center]
  {The row carries \emph{names}, never paths. That is what lets the same matrix
   drive two machines without a cell changing.};

\end{tikzpicture}
